\documentclass[12pt,a4paper]{article}

% ── Pacotes ──────────────────────────────────────────────────────
\usepackage[utf8]{inputenc}
\usepackage[T1]{fontenc}
\usepackage[brazil]{babel}
\usepackage[top=3cm,bottom=2cm,left=3cm,right=2cm,headheight=14pt]{geometry}
\usepackage{fancyhdr}
\usepackage{graphicx}
\usepackage{booktabs}
\usepackage{multirow}
\usepackage{array}
\usepackage{amsmath}
\usepackage{amssymb}
\usepackage{xcolor}
\usepackage{url}
\usepackage{hyperref}
\usepackage{caption}
\usepackage{subcaption}
\usepackage{setspace}
\usepackage{indentfirst}
\usepackage{parskip}

% ── Configurações ─────────────────────────────────────────────────
\hypersetup{
    colorlinks=true,
    linkcolor=blue!70!black,
    citecolor=blue!70!black,
    urlcolor=blue!70!black
}

\pagestyle{fancy}
\fancyhf{}
\fancyhead[L]{\small Lab.\ Experimentação de Software --- UFMG}
\fancyhead[R]{\small Lab02S02}
\fancyfoot[C]{\thepage}
\renewcommand{\headrulewidth}{0.4pt}

\setlength{\parindent}{1.5em}
\onehalfspacing

% ── Capa ──────────────────────────────────────────────────────────
\begin{document}

\begin{titlepage}
    \centering
    \vspace*{2cm}
    {\Large\bfseries Puc Minas Engenharia de Software\\[0.5em]
    Campus Praça da Liberdade\\[3em]}
    {\large\bfseries Lab02S02\\[0.8em]}
    {\LARGE\bfseries Um Estudo das Características de Qualidade\\[0.3em]
    de Sistemas Java no GitHub\\[3em]}
    {\large Integrantes: André Xavier Lazarini e João Pedro Guimarães\\[2em]}
    \vfill
    {\large Laboratório de Experimentação de Software\\[1em]
    2024/2025\\[2em]}
    \today
\end{titlepage}

% ── Sumário ───────────────────────────────────────────────────────
\tableofcontents
\newpage

% ═════════════════════════════════════════════════════════════════
\section{Introdução}
% ═════════════════════════════════════════════════════════════════

A qualidade do software é um aspecto fundamental no desenvolvimento de sistemas
robustos e de fácil manutenção. Métricas de software orientado a objetos,
como CBO (\emph{Coupling Between Objects}), DIT (\emph{Depth of Inheritance Tree})
e LCOM (\emph{Lack of Cohesion of Methods}), fornecem indicadores quantitativos
sobre o acoplamento, a herança e a coesão do código-fonte.

Neste trabalho, analisamos os \textbf{98 repositórios Java mais populares do GitHub}
(top por número de estrelas) com o objetivo de investigar a relação entre
características de processo --- popularidade, maturidade, atividade e tamanho ---
e as métricas de qualidade estrutural do código.

\subsection{Questões de Pesquisa}

\begin{description}
    \item[RQ01] Repositórios Java mais populares (número de estrelas) possuem
    código de maior qualidade (menor CBO, DIT e LCOM)?
    \item[RQ02] Repositórios Java mais maduros (tempo de existência) possuem
    código de maior qualidade?
    \item[RQ03] Repositórios Java mais ativos (número de releases) possuem
    código de maior qualidade?
    \item[RQ04] Repositórios Java maiores (LOC total) possuem código de maior
    qualidade?
\end{description}

% ═════════════════════════════════════════════════════════════════
\section{Metodologia}
% ═════════════════════════════════════════════════════════════════

\subsection{Coleta de Repositórios}

Os dados foram coletados em abril de 2025 por meio da GitHub Search API,
filtrando repositórios com a linguagem primária Java, ordenados por número
de estrelas em ordem decrescente
(\url{https://api.github.com/search/repositories?q=language:Java&sort=stars&per_page=100}).

Foram coletados os 98 repositórios do topo do ranking disponíveis no momento
da coleta. Para cada repositório foram registrados: nome completo, URL, número
de estrelas, forks, issues abertas, observadores, tamanho em KB, linguagem,
datas de criação/atualização/último push, branch padrão e total de releases
(obtido via GitHub Releases API usando paginação na cabeçalha \texttt{Link}).

\subsection{Análise Estática com CK}

Para cada repositório, realizou-se uma clonagem superficial
(\texttt{git clone --depth 1}) e executou-se a ferramenta
\textbf{CK v0.7.1}\footnote{\url{https://github.com/mauricioaniche/ck}}
conforme o comando:

\begin{center}
\texttt{java -jar ck.jar <repo\_dir> false 0 false <output\_dir>/}
\end{center}

A ferramenta gera um arquivo \texttt{class.csv} com métricas a nível de classe.
Para cada repositório foram calculadas as seguintes sumarizações das métricas:
mediana, média e desvio padrão de CBO, DIT e LCOM; mediana, média e total de LOC;
total de classes analisadas.

\subsection{Análise Estatística}

Dado que as distribuições das métricas não são normais (confirmado visualmente
pelas distribuições apresentadas na Seção~\ref{sec:bonus}), utilizou-se o
\textbf{coeficiente de correlação de Spearman} ($\rho$) para medir a associação
entre as métricas de processo e as métricas de qualidade.
Um valor-$p < 0,05$ foi adotado como critério de significância estatística.

% ═════════════════════════════════════════════════════════════════
\section{Resultados Gerais}
% ═════════════════════════════════════════════════════════════════

A Tabela~\ref{tab:central} apresenta as medidas de tendência central e dispersão
das principais métricas coletadas dos 98 repositórios analisados.

\begin{table}[htbp]
\centering
\caption{Medidas de tendência central e dispersão das métricas coletadas ($n = 98$).}
\label{tab:central}
\small
\begin{tabular}{lrrrrrr}
\toprule
\textbf{Métrica} & \textbf{Média} & \textbf{Mediana} & \textbf{D.P.} & \textbf{Mín.} & \textbf{Máx.} \\
\midrule
Estrelas           & 33.666  & 26.603  & 22.894    & 16.882  & 154.489   \\
Idade (anos)       & 9,70    & 9,35    & 3,28      & 2,30    & 16,14     \\
Total Releases     & 62,47   & 31,00   & 127,14    & 0       & 987       \\
CBO (mediana)      & 3,18    & 3,00    & 1,47      & 0,00    & 6,00      \\
DIT (mediana)      & 1,02    & 1,00    & 0,25      & 0,00    & 2,00      \\
LCOM (mediana)     & 0,63    & 0,00    & 1,72      & 0,00    & 13,00     \\
LOC (total)        & 95.126  & 37.277  & 155.658   & 0       & 772.133   \\
Nº Classes         & 2.092   & 672     & 3.788     & 0       & 24.751    \\
\bottomrule
\end{tabular}
\end{table}

Observa-se que a distribuição de estrelas e LOC é bastante assimétrica
(a média supera em muito a mediana), indicando a presença de repositórios
excepcionalmente grandes e populares. As métricas CBO e DIT apresentam valores
medianos baixos (3 e 1, respectivamente), enquanto LCOM tem mediana zero,
sugerindo alta coesão geral nos sistemas mais populares.

% ═════════════════════════════════════════════════════════════════
\section{RQ01 -- Popularidade e Qualidade}
% ═════════════════════════════════════════════════════════════════

\subsection{Motivação}

Repositórios com maior número de estrelas tendem a receber mais atenção da
comunidade, potencialmente levando a maior rigor nos padrões de codificação
e revisões de código. Investiga-se se essa atenção se reflete em métricas
de qualidade estrutural melhores.

\subsection{Hipótese}

$H_0$: Não há correlação entre a popularidade (estrelas) e as métricas de
qualidade (CBO, DIT, LCOM).\\
$H_1$: Repositórios mais populares apresentam menor CBO, DIT e LCOM.

\subsection{Resultados}

\begin{figure}[htbp]
    \centering
    \includegraphics[width=\textwidth]{figs/rq01_scatter.png}
    \caption{Dispersão entre Popularidade (estrelas) e métricas de qualidade.
    $\rho$ = coeficiente de Spearman; * $p < 0,05$.}
    \label{fig:rq01}
\end{figure}

\begin{table}[htbp]
\centering
\caption{Correlação de Spearman: Popularidade $\times$ Qualidade ($n = 98$).}
\label{tab:rq01_spearman}
\begin{tabular}{lrrr}
\toprule
\textbf{Qualidade} & \textbf{$\rho$} & \textbf{valor-$p$} & \textbf{Significativo?} \\
\midrule
CBO (mediana) & $-0{,}134$ & $0{,}189$ & Não \\
DIT (mediana) & $-0{,}063$ & $0{,}535$ & Não \\
LCOM (mediana)& $-0{,}117$ & $0{,}251$ & Não \\
\bottomrule
\end{tabular}
\end{table}

\begin{table}[htbp]
\centering
\caption{Medidas de qualidade por quartil de popularidade (estrelas).}
\label{tab:rq01_quartis}
\small
\begin{tabular}{llrrr}
\toprule
\textbf{Quartil} & \textbf{Métrica} & \textbf{$n$} & \textbf{Média} & \textbf{Mediana} \\
\midrule
\multirow{3}{*}{Q1 (baixo)} & CBO & 25 & 3,48 & 3,00 \\
                             & DIT & 25 & 1,00 & 1,00 \\
                             & LCOM& 25 & 0,96 & 0,00 \\
\midrule
\multirow{3}{*}{Q2}          & CBO & 24 & 2,92 & 3,00 \\
                             & DIT & 24 & 1,04 & 1,00 \\
                             & LCOM& 24 & 0,33 & 0,00 \\
\midrule
\multirow{3}{*}{Q3}          & CBO & 24 & 3,75 & 4,00 \\
                             & DIT & 24 & 1,08 & 1,00 \\
                             & LCOM& 24 & 0,54 & 0,00 \\
\midrule
\multirow{3}{*}{Q4 (alto)}   & CBO & 25 & 2,60 & 3,00 \\
                             & DIT & 25 & 0,96 & 1,00 \\
                             & LCOM& 25 & 0,68 & 0,00 \\
\bottomrule
\end{tabular}
\end{table}

\subsection{Discussão}

Nenhuma das correlações entre popularidade e qualidade foi estatisticamente
significativa (Tabela~\ref{tab:rq01_spearman}). Os coeficientes negativos
para CBO ($\rho = -0{,}134$) e LCOM ($\rho = -0{,}117$) sugerem uma tendência
muito tênue de redução do acoplamento e melhoria da coesão com o aumento da
popularidade, mas essa tendência não passa do limiar de significância.

A análise por quartis (Tabela~\ref{tab:rq01_quartis}) também não revela
padrão monotônico claro. O Q4 (mais populares) apresenta CBO médio de 2,60,
ligeiramente inferior ao Q1 (3,48), mas as variações entre quartis são pequenas.

\textbf{Conclusão RQ01}: Não há evidência de que repositórios mais populares
no GitHub possuam código Java de maior qualidade estrutural.

% ═════════════════════════════════════════════════════════════════
\section{RQ02 -- Maturidade e Qualidade}
% ═════════════════════════════════════════════════════════════════

\subsection{Motivação}

Repositórios mais antigos passaram por mais ciclos de refatoração, revisão
e evolução. A hipótese é que a maturidade temporal esteja associada a
arquiteturas mais estáveis e, portanto, métricas de qualidade melhores.

\subsection{Resultados}

\begin{figure}[htbp]
    \centering
    \includegraphics[width=\textwidth]{figs/rq02_scatter.png}
    \caption{Dispersão entre Maturidade (anos) e métricas de qualidade.
    $\rho$ = coeficiente de Spearman; ** $p < 0,01$.}
    \label{fig:rq02}
\end{figure}

\begin{table}[htbp]
\centering
\caption{Correlação de Spearman: Maturidade $\times$ Qualidade ($n = 98$).}
\label{tab:rq02_spearman}
\begin{tabular}{lrrr}
\toprule
\textbf{Qualidade} & \textbf{$\rho$} & \textbf{valor-$p$} & \textbf{Significativo?} \\
\midrule
CBO (mediana) & $+0{,}157$ & $0{,}123$ & Não \\
DIT (mediana) & $+0{,}318$ & $0{,}001$ & \textbf{Sim} (**) \\
LCOM (mediana)& $-0{,}011$ & $0{,}911$ & Não \\
\bottomrule
\end{tabular}
\end{table}

\begin{table}[htbp]
\centering
\caption{Medidas de qualidade por quartil de maturidade (idade em anos).}
\label{tab:rq02_quartis}
\small
\begin{tabular}{llrrr}
\toprule
\textbf{Quartil} & \textbf{Métrica} & \textbf{$n$} & \textbf{Média} & \textbf{Mediana} \\
\midrule
\multirow{3}{*}{Q1 (mais novos)} & CBO & 25 & 2,68 & 2,00 \\
                                  & DIT & 25 & 0,96 & 1,00 \\
                                  & LCOM& 25 & 0,48 & 0,00 \\
\midrule
\multirow{3}{*}{Q2}               & CBO & 24 & 3,17 & 3,00 \\
                                  & DIT & 24 & 0,96 & 1,00 \\
                                  & LCOM& 24 & 1,44 & 0,00 \\
\midrule
\multirow{3}{*}{Q3}               & CBO & 24 & 3,67 & 4,00 \\
                                  & DIT & 24 & 1,00 & 1,00 \\
                                  & LCOM& 24 & 0,17 & 0,00 \\
\midrule
\multirow{3}{*}{Q4 (mais velhos)} & CBO & 25 & 3,24 & 3,00 \\
                                  & DIT & 25 & 1,16 & 1,00 \\
                                  & LCOM& 25 & 0,46 & 0,00 \\
\bottomrule
\end{tabular}
\end{table}

\subsection{Discussão}

A única correlação significativa encontrada foi entre maturidade e DIT
($\rho = +0{,}318$, $p = 0{,}001$). Curiosamente, a correlação é \emph{positiva},
indicando que repositórios mais antigos tendem a apresentar hierarquias de
herança \emph{mais profundas}, e não mais rasas.

Uma explicação plausível é que sistemas mais maduros acumulam mais camadas
de abstração ao longo do tempo, refletindo decisões de design que favorecem
extensibilidade via herança. Isso não necessariamente implica pior qualidade,
mas sim uma escolha arquitetural mais frequente em sistemas legados Java.

CBO e LCOM não apresentaram correlações significativas com a maturidade,
sugerindo que o tempo de existência isoladamente não determina acoplamento
ou coesão.

\textbf{Conclusão RQ02}: Repositórios mais maduros tendem a ter hierarquias
de herança ligeiramente mais profundas (DIT), mas não diferem significativamente
em CBO e LCOM. A hipótese de que maturidade implica código de maior qualidade
não é suportada pelos dados.

% ═════════════════════════════════════════════════════════════════
\section{RQ03 -- Atividade e Qualidade}
% ═════════════════════════════════════════════════════════════════

\subsection{Motivação}

O número de releases reflete a frequência de entregas formais de software.
Projetos com muitas releases tendem a ter processos de engenharia mais maduros,
potencialmente incluindo práticas de refatoração contínua.

\subsection{Resultados}

\begin{figure}[htbp]
    \centering
    \includegraphics[width=\textwidth]{figs/rq03_scatter.png}
    \caption{Dispersão entre Atividade (releases) e métricas de qualidade.
    $\rho$ = coeficiente de Spearman; * $p < 0,05$; *** $p < 0,001$.}
    \label{fig:rq03}
\end{figure}

\begin{table}[htbp]
\centering
\caption{Correlação de Spearman: Atividade $\times$ Qualidade ($n = 98$).}
\label{tab:rq03_spearman}
\begin{tabular}{lrrr}
\toprule
\textbf{Qualidade} & \textbf{$\rho$} & \textbf{valor-$p$} & \textbf{Significativo?} \\
\midrule
CBO (mediana) & $+0{,}379$ & $< 0{,}001$ & \textbf{Sim} (***) \\
DIT (mediana) & $+0{,}201$ & $0{,}048$   & \textbf{Sim} (*) \\
LCOM (mediana)& $+0{,}107$ & $0{,}294$   & Não \\
\bottomrule
\end{tabular}
\end{table}

\begin{table}[htbp]
\centering
\caption{Medidas de qualidade por quartil de atividade (total releases).
Nota: Q1 agrupou repositórios com 0 releases (50 repositórios).}
\label{tab:rq03_quartis}
\small
\begin{tabular}{llrrr}
\toprule
\textbf{Quartil} & \textbf{Métrica} & \textbf{$n$} & \textbf{Média} & \textbf{Mediana} \\
\midrule
\multirow{3}{*}{Q1 (menos ativos)} & CBO & 50 & 2,64 & 3,00 \\
                                    & DIT & 50 & 0,98 & 1,00 \\
                                    & LCOM& 50 & 0,69 & 0,00 \\
\midrule
\multirow{3}{*}{Q2}                 & CBO & 23 & 4,00 & 4,00 \\
                                    & DIT & 23 & 1,09 & 1,00 \\
                                    & LCOM& 23 & 0,57 & 0,00 \\
\midrule
\multirow{3}{*}{Q3 (mais ativos)}   & CBO & 25 & 3,52 & 4,00 \\
                                    & DIT & 25 & 1,04 & 1,00 \\
                                    & LCOM& 25 & 0,58 & 0,00 \\
\bottomrule
\end{tabular}
\end{table}

\subsection{Discussão}

Encontrou-se correlação significativa entre atividade e CBO ($\rho = +0{,}379$,
$p < 0{,}001$) e DIT ($\rho = +0{,}201$, $p = 0{,}048$). Ambas as correlações
são \emph{positivas}, contrariando a hipótese inicial: repositórios com mais
releases tendem a ter \emph{maior} acoplamento e herança mais profunda,
não menor.

Uma interpretação é que projetos com mais releases são sistemas mais complexos
e maduros (frameworks, bibliotecas, servidores), que naturalmente apresentam
maior acoplamento entre componentes. Nesse contexto, o CBO elevado pode refletir
a riqueza de integrações do sistema, não necessariamente negligência com
qualidade.

LCOM não apresentou correlação significativa, indicando que a atividade
de releases não influencia a coesão interna das classes.

\textbf{Conclusão RQ03}: Repositórios mais ativos (com mais releases) tendem
a apresentar maior acoplamento (CBO) e hierarquias levemente mais profundas
(DIT). Isso sugere que projetos mais ativos são intrinsecamente mais complexos,
o que se reflete nessas métricas.

% ═════════════════════════════════════════════════════════════════
\section{RQ04 -- Tamanho e Qualidade}
% ═════════════════════════════════════════════════════════════════

\subsection{Motivação}

Sistemas maiores, medidos em linhas de código, tendem a ser mais complexos.
Investiga-se se esse tamanho se reflete em maiores valores de acoplamento,
profundidade de herança e falta de coesão.

\subsection{Resultados}

\begin{figure}[htbp]
    \centering
    \includegraphics[width=\textwidth]{figs/rq04_scatter.png}
    \caption{Dispersão entre Tamanho (LOC total) e métricas de qualidade.
    $\rho$ = coeficiente de Spearman; * $p < 0,05$; *** $p < 0,001$.}
    \label{fig:rq04}
\end{figure}

\begin{table}[htbp]
\centering
\caption{Correlação de Spearman: Tamanho $\times$ Qualidade ($n = 98$).}
\label{tab:rq04_spearman}
\begin{tabular}{lrrr}
\toprule
\textbf{Qualidade} & \textbf{$\rho$} & \textbf{valor-$p$} & \textbf{Significativo?} \\
\midrule
CBO (mediana) & $+0{,}348$ & $< 0{,}001$ & \textbf{Sim} (***) \\
DIT (mediana) & $+0{,}234$ & $0{,}020$   & \textbf{Sim} (*) \\
LCOM (mediana)& $+0{,}108$ & $0{,}291$   & Não \\
\bottomrule
\end{tabular}
\end{table}

\begin{table}[htbp]
\centering
\caption{Medidas de qualidade por quartil de tamanho (LOC total).}
\label{tab:rq04_quartis}
\small
\begin{tabular}{llrrr}
\toprule
\textbf{Quartil} & \textbf{Métrica} & \textbf{$n$} & \textbf{Média} & \textbf{Mediana} \\
\midrule
\multirow{3}{*}{Q1 (menores)} & CBO & 25 & 2,08 & 2,00 \\
                               & DIT & 25 & 0,92 & 1,00 \\
                               & LCOM& 25 & 0,76 & 0,00 \\
\midrule
\multirow{3}{*}{Q2}            & CBO & 24 & 3,46 & 3,00 \\
                               & DIT & 24 & 1,04 & 1,00 \\
                               & LCOM& 24 & 0,58 & 0,00 \\
\midrule
\multirow{3}{*}{Q3}            & CBO & 24 & 3,58 & 3,00 \\
                               & DIT & 24 & 1,04 & 1,00 \\
                               & LCOM& 24 & 0,21 & 0,00 \\
\midrule
\multirow{3}{*}{Q4 (maiores)}  & CBO & 25 & 3,64 & 4,00 \\
                               & DIT & 25 & 1,08 & 1,00 \\
                               & LCOM& 25 & 0,96 & 0,00 \\
\bottomrule
\end{tabular}
\end{table}

\subsection{Discussão}

Assim como RQ03, encontrou-se correlação positiva e significativa entre tamanho
e CBO ($\rho = +0{,}348$, $p < 0{,}001$) e entre tamanho e DIT ($\rho = +0{,}234$,
$p = 0{,}020$). A análise por quartis confirma o padrão: o CBO médio sobe de
2,08 no Q1 para 3,64 no Q4.

Essa relação é esperada: sistemas maiores contêm mais classes que necessariamente
interagem entre si, elevando o CBO. O aumento de DIT também é natural em
sistemas grandes, que tipicamente utilizam mais frameworks e possuem hierarquias
de abstração mais elaboradas.

LCOM não apresentou correlação significativa com o tamanho, o que é interessante:
mesmo sistemas grandes mantêm coesão interna mediana de zero, sugerindo que
as classes individuais são focadas em suas responsabilidades.

\textbf{Conclusão RQ04}: Sistemas Java maiores (em LOC) apresentam maior
acoplamento (CBO) e hierarquias mais profundas (DIT), o que é esperado dado
o aumento de complexidade. Coesão (LCOM) não é afetada pelo tamanho.

% ═════════════════════════════════════════════════════════════════
\section{Bônus -- Análise Complementar}
\label{sec:bonus}
% ═════════════════════════════════════════════════════════════════

\subsection{Distribuição das Métricas de Qualidade}

\begin{figure}[htbp]
    \centering
    \includegraphics[width=\textwidth]{figs/distribuicoes_qualidade.png}
    \caption{Distribuição de CBO, DIT e LCOM nos 98 repositórios analisados.
    Linha tracejada: mediana. Linha pontilhada vermelha: média.}
    \label{fig:distribuicoes}
\end{figure}

A Figura~\ref{fig:distribuicoes} mostra que todas as três métricas de qualidade
apresentam distribuições assimétricas positivas. CBO concentra-se entre 0 e 6,
com mediana 3. DIT tem distribuição muito concentrada em torno de 1, indicando
que a maioria dos projetos usa herança pouco profunda. LCOM apresenta muitos
zeros (alta coesão) mas com cauda longa, refletindo algumas classes com baixa coesão.

\subsection{Boxplots por Quartis de Processo}

\begin{figure}[htbp]
    \centering
    \includegraphics[width=\textwidth]{figs/boxplots_quartis.png}
    \caption{Distribuição das métricas de qualidade por quartis de cada métrica
    de processo. Cada linha corresponde a uma RQ; cada coluna a uma métrica
    de qualidade.}
    \label{fig:boxplots}
\end{figure}

Os boxplots da Figura~\ref{fig:boxplots} permitem visualizar a variação das
métricas de qualidade dentro de cada quartil. Observa-se que o CBO tem maior
variabilidade nos quartis superiores de releases e LOC total, consistente
com as correlações positivas encontradas.

\subsection{Matriz de Correlação de Spearman}

\begin{figure}[htbp]
    \centering
    \includegraphics[width=0.85\textwidth]{figs/heatmap_correlacoes.png}
    \caption{Heatmap da matriz de correlação de Spearman entre todas as métricas
    coletadas. Valores positivos (verde) indicam correlação direta; negativos
    (vermelho) indicam correlação inversa.}
    \label{fig:heatmap}
\end{figure}

O heatmap da Figura~\ref{fig:heatmap} revela padrões interessantes além das
quatro RQs:
\begin{itemize}
    \item LOC total e número de classes apresentam forte correlação entre si,
    como esperado.
    \item Releases correlaciona positivamente com CBO e LOC, sugerindo que
    projetos ativos são intrinsecamente maiores e mais acoplados.
    \item DIT correlaciona fracamente com a maioria das métricas, exceto com
    maturidade (RQ02) e releases/LOC (RQ03/RQ04).
\end{itemize}

\subsection{Resumo Consolidado das Correlações}

\begin{table}[htbp]
\centering
\caption{Resumo das correlações de Spearman para todas as questões de pesquisa.}
\label{tab:resumo}
\small
\begin{tabular}{llrrrc}
\toprule
\textbf{RQ} & \textbf{Processo} & \textbf{Qualidade} & \textbf{$\rho$} & \textbf{valor-$p$} & \textbf{Sig.} \\
\midrule
RQ01 & Estrelas & CBO  & $-0{,}134$ & $0{,}189$ & \\
RQ01 & Estrelas & DIT  & $-0{,}063$ & $0{,}535$ & \\
RQ01 & Estrelas & LCOM & $-0{,}117$ & $0{,}251$ & \\
\midrule
RQ02 & Idade (anos) & CBO  & $+0{,}157$ & $0{,}123$ & \\
RQ02 & Idade (anos) & DIT  & $+0{,}318$ & $0{,}001$ & ** \\
RQ02 & Idade (anos) & LCOM & $-0{,}011$ & $0{,}911$ & \\
\midrule
RQ03 & Releases & CBO  & $+0{,}379$ & $< 0{,}001$ & *** \\
RQ03 & Releases & DIT  & $+0{,}201$ & $0{,}048$   & * \\
RQ03 & Releases & LCOM & $+0{,}107$ & $0{,}294$   & \\
\midrule
RQ04 & LOC total & CBO  & $+0{,}348$ & $< 0{,}001$ & *** \\
RQ04 & LOC total & DIT  & $+0{,}234$ & $0{,}020$   & * \\
RQ04 & LOC total & LCOM & $+0{,}108$ & $0{,}291$   & \\
\bottomrule
\multicolumn{6}{l}{\small * $p < 0{,}05$; \quad ** $p < 0{,}01$; \quad *** $p < 0{,}001$}
\end{tabular}
\end{table}

% ═════════════════════════════════════════════════════════════════
\section{Conclusão}
% ═════════════════════════════════════════════════════════════════

Este trabalho analisou 98 dos repositórios Java mais populares do GitHub,
aplicando a ferramenta CK para extração de métricas de qualidade estrutural
(CBO, DIT, LCOM) e investigando suas relações com características de processo
(popularidade, maturidade, atividade e tamanho).

Os principais achados são:

\begin{enumerate}
    \item \textbf{Popularidade (RQ01)}: Não há correlação significativa entre
    o número de estrelas e as métricas de qualidade analisadas. A hipótese
    de que projetos mais populares são de maior qualidade estrutural não
    foi confirmada.

    \item \textbf{Maturidade (RQ02)}: Repositórios mais antigos tendem a ter
    hierarquias de herança ligeiramente mais profundas (DIT, $\rho = +0{,}318$),
    possivelmente devido ao acúmulo de camadas de abstração ao longo do tempo.

    \item \textbf{Atividade (RQ03)}: Repositórios com mais releases apresentam
    maior acoplamento (CBO, $\rho = +0{,}379$) e herança levemente mais profunda
    (DIT, $\rho = +0{,}201$), refletindo a complexidade inerente de sistemas
    com ciclos de release mais frequentes.

    \item \textbf{Tamanho (RQ04)}: Sistemas maiores (LOC) têm maior acoplamento
    (CBO, $\rho = +0{,}348$) e hierarquias mais profundas (DIT, $\rho = +0{,}234$),
    o que é esperado dado o crescimento da complexidade estrutural.
\end{enumerate}

Um achado transversal relevante é que o LCOM não apresentou correlação
significativa com nenhuma das métricas de processo, sugerindo que a coesão
interna das classes é uma propriedade relativamente estável e independente
de popularidade, maturidade, atividade e tamanho dos projetos.

\subsection{Ameaças à Validade}

\begin{itemize}
    \item \textbf{Validade interna}: A coleta foi feita em um momento específico;
    projetos em evolução ativa podem ter suas métricas alteradas rapidamente.

    \item \textbf{Validade externa}: A amostra é restrita aos 98 projetos
    mais populares. Projetos menos famosos podem exibir padrões diferentes.

    \item \textbf{Construção}: CBO, DIT e LCOM são métricas de acoplamento
    estrutural; qualidade percebida por desenvolvedores pode não se correlacionar
    com essas métricas estáticas.

    \item \textbf{Tamanho da amostra}: Com 98 repositórios (de um total de 1000
    planejados), algumas correlações podem não ter poder estatístico suficiente
    para se revelar.
\end{itemize}

\vspace{1em}
\noindent
\textbf{Dados e código disponíveis em}:
\url{https://github.com/andrexlazarini/LabExperimentacao-2}

\end{document}
